\begin{abstract}
Local-context world--action models (WAMs) couple visual prediction with control but leave task-critical earlier interactions outside their decision context.
Memory-augmented embodied policies extend access to history, yet retaining past information does not establish whether a recalled action remains appropriate.
The central challenge is to use ongoing interaction as evidence for deciding which past behavior to reuse and how it should shape action generation.
We introduce \method{} \textit{\textbf{(\textbf{W}orld--\textbf{A}ction--\textbf{R}etrospection \textbf{M}odel)}}, organized around \emph{event-grounded retrospection}: current context and \emph{situated recollection} ground inference of a candidate-independent transition requirement, while a cross-episode \emph{sensorimotor repertoire} supplies state--action--effect events.
Adapted repertoire actions are assessed by comparing their predicted effects with this learned forecast; \emph{consequence-gated source transport} then jointly regulates proposal-informed stochastic initialization and candidate-derived conditioning, turning recalled behavior into a revisable prior for flow-based action generation.
\method{} achieves 98.7\% success on LIBERO and 83.9\% on RMBench with task-specialist models, including 87.3\% on tasks requiring multiple past observations.
Component comparisons show complementary benefits from recollection and repertoire; complete-event distraction tests show less degradation with explicit consequence comparison.
These findings support using current evidence to govern both the selection of past behavior and its influence on action generation.
\end{abstract}
